\expandafter\def\csname HistoricalNoteText8\endcsname{When I had the good fortune some time ago to meet Professor \textsc{Gauss} in person, I saw among his papers the essay that follows, drafted several years earlier and still unpublished. It contains the author’s earlier method of determining planetary orbits. On a quick reading of this overview I soon became convinced that the method the author develops here for obtaining an initial approximation to two distances of the planet from the Earth differs substantially from the one he has now presented publicly in his larger work. I therefore asked his permission to publish the essay, assuming that anyone familiar with the subject would find it interesting to know the different routes by which the author succeeded in reaching the complete solution that we outlined for our readers in the preceding issue. At first I intended to accompany the essay with some remarks comparing the author’s earlier and later methods. But if those remarks were really to clarify the matter, they would have to be somewhat lengthy and would still remain unclear without reference to the work itself. It therefore seemed better to give the astronomical readers of this journal the whole essay, which is intended chiefly for readers familiar with the subject who have the work itself at hand, exactly as the author wrote it six years ago, without further additions.\hfill\textsc{von Lindenau.}}
\clearpage\SourceLeaf{153}%
\pdfbookmark[1]{An Outline of the Methods Used to Determine the Orbits of the Two New Primary Planets}{summarische1809}
\fancyhead[L]{\small OUTLINE OF METHODS FOR DETERMINING ORBITS}
\begin{center}
{\small AN OUTLINE OF THE METHODS USED TO DETERMINE THE ORBITS OF THE TWO NEW PRIMARY PLANETS\HistoricalNote{8}.\par}
\medskip\hrule\medskip
{\small Monatliche Correspondenz zur Beförderung der Erd- und Himmels-Kunde,\par
edited by Baron \textsc{von Zach}. September 1809.\par}
\medskip\hrule
\end{center}\medskip
\begin{center}1.\end{center}
Determining the orbit of a celestial body from a short series of observations without assuming a circle or a parabola rests on two requirements. I. One must have a means of finding the orbit that fits three given complete observations. II. One must be able to refine the orbit thus found so that the differences between the calculated positions and the full set of observations are as small as possible.

\SourceLeaf{154}%
The easiest way to meet the II\textsuperscript{nd} requirement seems to be to reduce it to the I\textsuperscript{st}. At times $t$, $t'$, $t''$, and so on, let the observed positions be $m$, $m'$, $m''$ (each of which will have two components). Let the positions calculated from the known elements $e$ (which are to be regarded as six in number) be $p$, $p'$, $p''$, and so on. Finally, let the positions calculated from the elements $f$ (which are still to be regarded as undetermined) be $q$, $q'$, $q''$, and so on. The differences for the elements $e$ are thus
\[
p-m,\quad p'-m',\quad p''-m''\ \text{and so on}
\]
For the elements $f$, however, the differences are
\[
(p-m)+(q-p),\quad(p'-m')+(q'-p'),\quad(p''-m'')+(q''-p'')\ \text{and so on}
\]
The latter are to be made as small as possible, with no regular pattern left.

The differences $q-p$, $q'-p'$, and so on are functions of time insofar as the elements $f$ are treated as constant. Since by the nature of the problem these differences will themselves be small, given the short span of the observations one may assume that, if the differences are given for two outer times and one middle time, interpolation will give them accurately enough for the intervening times. Denote them at those three times by $x$, $y$, $z$ (each with two components). By familiar principles of interpolation theory they will then have the linear form $\alpha x+\beta y+\gamma z$, where the coefficients $\alpha$, $\beta$, $\gamma$ depend on time. Thus the differences for the elements $f$ will have the following form at times $t$, $t'$, $t''$, and so on.
\[
\begin{gathered}
p-m+\alpha x+\beta y+\gamma z\\
p'-m'+\alpha'x+\beta'y+\gamma'z\\
p''-m''+\alpha''x+\beta''y+\gamma''z\ \text{and so on}
\end{gathered}
\]
Here everything except $x$, $y$, $z$ is known. It will then be easy to judge which values of $x$, $y$, $z$ are most suitable. \emph{A fully systematic procedure can indeed be given for calculating these values.} But a certain feel for the problem will always guide one just as surely.

Since the II\textsuperscript{nd} requirement is evidently reduced to the first as soon as $x$, $y$, $z$ have been determined, we can restrict ourselves to the first.

\SourceLeaf{155}%
\begin{center}\emph{Determining the orbit from three complete observations.}\end{center}
\begin{center}2.\end{center}
It would be easy enough to express the relation of the six unknown quantities to the given ones in six equations. These equations, however, would be much too unwieldy to be of any use. One must therefore be content to proceed in stages to the orbit that reproduces the three observations exactly. Clearly every method that serves this purpose must eventually give the same result. The quality of the final result therefore measures the precision of the observations used. It gives no measure of the merit of the method. A method can be judged only by how many stages it takes and how easily they can be carried out. If a method could not always reproduce the three observations exactly whenever desired, then calling it a poorer method would already concede too much. It would be no method at all.

The investigation thus falls into two parts, a \emph{first approximation} and the \emph{methods of correction}. The former will rest on certain nearly exact relations drawn from the nature of the problem. They are increasingly accurate the closer together the observations are and, mathematically speaking, are strictly exact for observations infinitely close together. Thus, the closer the observations on which one bases the calculation, the less the departure of these relations from the truth affects it. This makes the correction methods easier to use or even makes them unnecessary. One must bear in mind, however, that when the observations are close together, small errors in them have a very strong and sometimes enormous effect on the elements. The subsequent refinement using the whole series of observations, which was called the second requirement above, therefore becomes all the more difficult. It is hard to give general rules for choosing the most suitable observations. Their quality, whether their positions are more or less favourable, and the nature of the orbit itself all have to be taken into account. For Ceres the outermost observations, 41 days apart, could be used successfully right away for the first approximation. The correction calculations were no trouble at all. In calculating the II\textsuperscript{nd} orbit of Pallas, the previous approximations were likewise left aside, and the first
\SourceLeaf{156}%
approximation method was applied afresh straight away to the observations spanning 27 days. For orbits closer to a parabola, in which the geocentric motion is very fast, it would probably be preferable to start with somewhat shorter intervals. Judgment sharpened by experience is the best guide here.

\begin{center}
\emph{Main stages of the first approximation.}\par\medskip
First stage.\par\smallskip
\emph{Finding approximate distances from the Earth at the two outer observations.}
\end{center}
\begin{center}3.\end{center}
To make it easier to keep track of the many symbols needed in the following investigation, the same letters will be used for corresponding quantities of the Earth $P$ and the observed planet $p$ throughout, with capitals for the Earth and lowercase letters for the planet. If a letter occurs unprimed as well as with one and two primes, the second and third forms are to be understood as referring to the second and third observations at times $\tau'$, $\tau''$ in the same way as the first form refers to the observation at time $\tau$. The time $\tau'$ need not in itself fall between $\tau$ and $\tau''$. It is most advantageous, however, to use the following rules when $\tau'$ is approximately midway between $\tau$ and $\tau''$.

$S$ is the position of the Sun (in space), regarded as fixed. $p$ is the position of the planet $p$ at time $\tau$. This explains $p'$, $p''$, $P$, $P'$, $P''$.

$\mathfrak{X}$, $\mathfrak{Y}$, $\mathfrak{Z}$ are three arbitrary fixed planes that intersect one another at right angles at the centre of the Sun.

$x$, $y$, $z$ are the perpendicular distances of the planet $p$ from these three planes at time $\tau$. Hence $x'$, $y'$, $z'$; $x''$, $y''$, $z''$; $X$, $Y$, $Z$; $X'$, $Y'$, $Z'$; $X''$, $Y''$, $Z''$.
\[
\left.\begin{aligned}
\xi&=x-X\\
\eta&=y-Y\\
\zeta&=z-Z
\end{aligned}\right|\quad\text{Hence }\xi',\ \eta',\ \zeta';\ \xi'',\ \eta'',\ \zeta''
\]

Thus $\xi$, $\eta$, $\zeta$ are the perpendicular distances of the planet $p$ from three moving planes drawn parallel to $\mathfrak{X}$, $\mathfrak{Y}$, $\mathfrak{Z}$ through $P$.

\SourceLeaf{157}%
\[
\begin{array}{c|c|c|c|c|l}
r&&\text{of }p&&S&\text{all positive.}\\
R&\text{distances}&\text{of }P&\text{from}&S&\text{Hence }r'\ \text{and so on}\\
\rho&&\text{of }p&&P&
\end{array}
\]
\[
\begin{array}{c|c|c|c|l}
b&&Sp&&\mathfrak{Z}\\
B&\text{angle of the line}&SP&\text{with the plane}&\mathfrak{Z}\\
\beta&&Pp&&\text{parallel to }\mathfrak{Z}
\end{array}
\]
\[
\left.\begin{aligned}
d&=r\cos b\\D&=R\cos B\\\delta&=\rho\cos\beta
\end{aligned}\right|\quad
\begin{gathered}\text{that is, distances projected onto the plane }\mathfrak{Z}\\\text{and the plane parallel to it}\end{gathered}
\]
\[
\begin{array}{c|c|l}
l&&\mathfrak{Y}\\
L&\text{angle of this projection with the plane}&\mathfrak{Y}\\
\lambda&&\text{parallel to }\mathfrak{Y}\text{.}
\end{array}
\]
The angles $b$ and $l$ are to be taken as positive on the same sides of $\mathfrak{Z}$ and $\mathfrak{Y}$ on which $z$ and $y$ are taken as positive. One can always keep the angles $b$ between the limits $-90^0$, $+90^0$ (so that the $d$ and so on always remain positive). The angles $l$, on the other hand, can be allowed to increase from $0$ to $360^0$, taking them as $\left\{\begin{smallmatrix}0\\180^0\end{smallmatrix}\right\}$ where the $x$ are $\left\{\begin{smallmatrix}\text{positive}\\\text{negative}\end{smallmatrix}\right\}$. In this way one obtains
\[
\left.\begin{aligned}
x&=r\cos b\cos l=d\cos l\\
y&=r\cos b\sin l=d\sin l\\
z&=r\sin b\phantom{\cos l}=d\operatorname{tang}b
\end{aligned}\right|\quad
\begin{gathered}\text{and similar equations for}\\x'\ \text{and so on}\quad X\ \text{and so on}\quad\xi\ \text{and so on}\end{gathered}
\]

We take the orbits of $p$ and $P$ to lie in planes, disregarding the external influences that affect them. We put the longitude of $p$ in its orbit at time $\tau$ $=v$ (hence $v'$, $v''$; $V$, $V'$, $V''$) and put $\frac12r'r''\sin(v''-v')=f$, $\frac12r''r\sin(v-v'')=f'$, $\frac12rr'\sin(v'-v)=f''$. Accordingly $f$, $-f'$, $f''$ are the positive areas of the triangles $p'Sp''$, $pSp''$, $pSp'$ (assuming that $p$ moves prograde and that $\tau'$ lies between $\tau$ and $\tau''$. The signs are easy to assign in other cases). Similarly for $F$, $F'$, $F''$. By $g$, $-g'$, $g''$, $G$, $-G'$, $G''$ we denote the areas of the sectors of the full orbit corresponding to these triangles, and take their signs to agree with those of $f$, $-f'$, $f''$, $F$, $-F'$, $F''$. Thus $g$, $g'$, $g''$ and $G$, $G'$, $G''$ are proportional to the time intervals $\tau''-\tau'$, $\tau-\tau''$, $\tau'-\tau$.

\SourceLeaf{158}%
The orbits of $p$ and $P$ are conic sections, and we denote their semimajor axes by $a$, $A$. We set the eccentricity of the orbit of $p$ $=e=\sin\varphi$ (for an ellipse), so that $a\cos\varphi^2=k$ is the half parameter. Longitude of the aphelion of $p$ in its orbit $=\pi$. Mean longitude $=m$ (hence $m'$, $m''$, $M$, $M'$, $M''$). Other symbols will be introduced in the course of the investigation.

\begin{center}4.\end{center}
Since $p$, $p'$, $p''$ and $S$ lie in one plane, a familiar theorem gives
\[
0=xy'z''+x'y''z+x''yz'-xy''z'-x'yz''-x''y'z
\]
and from this it follows that of the following nine quantities
\[
\begin{array}{ccc}
y'z''-y''z'&y''z-yz''&yz'-y'z\\
z'x''-z''x'&z''x-zx''&zx'-z'x\\
x'y''-x''y'&x''y-xy''&xy'-x'y
\end{array}
\]
the upper three are proportional, respectively, to the middle three and to the lower three. It is easy to see the following.

I. The quantities $f$, $f'$, $f''$ are also proportional to these same quantities. The upper, middle and lower sets of three are simply twice the areas of the projections of the triangles whose areas are $f$, $f'$, $f''$ onto the fundamental planes $\mathfrak{X}$, $\mathfrak{Y}$, $\mathfrak{Z}$. In each case, twice the projected areas thus stand to the original areas as twice the cosine of the inclination of the orbit of $p$ to the corresponding plane stands to unity. (A full treatment would still need some remarks here about the signs, but these can easily be dispensed with by direct calculation.)

II. If one multiplies the upper three, the middle three or the lower three by $x$, $x'$, $x''$, by $y$, $y'$, $y''$ or by $z$, $z'$, $z''$, respectively, the sum of the products becomes $=0$. From this one readily concludes
\[
\begin{array}{r|c|l}
0=fx+f'x'+f''x''&\text{By exactly}&0=FX+F'X'+F''X''\\
0=fy+f'y'+f''y''&\text{analogous reasoning}&0=FY+F'Y'+F''Y''\\
0=fz+f'z'+f''z''&\text{one obtains}&0=FZ+F'Z'+F''Z''
\end{array}
\]
\SourceLeaf{159}%
From these one can readily derive the following three equations
\[
\begin{aligned}
&(F+F'')(f\xi+f'\xi'+f''\xi'')\\
&\qquad=(Ff''-F''f)(X-X'')+[F'(f+f'')-(F+F'')f']X'\\
&(F+F'')(f\eta+f'\eta'+f''\eta'')\\
&\qquad=(Ff''-F''f)(Y-Y'')+[F'(f+f'')-(F+F'')f']Y'\\
&(F+F'')(f\zeta+f'\zeta'+f''\zeta'')\\
&\qquad=(Ff''-F''f)(Z-Z'')+[F'(f+f'')-(F+F'')f']Z'
\end{aligned}
\]

From these three equations we derive four others as follows
\[
\begin{array}{c|c|c|c|l}
\text{first by}&\text{then by}&\text{then by}&\text{finally by}&\text{multiply}\\
\eta\zeta''-\eta''\zeta&\eta Z'-Y'\zeta&\eta'Z'-Y'\zeta'&\eta''Z'-Y'\zeta''&\text{and}\\
\zeta\xi''-\zeta''\xi&\zeta X'-Z'\xi&\zeta'X'-Z'\xi'&\zeta''X'-Z'\xi''&\text{add the}\\
\xi\eta''-\xi''\eta&\xi Y'-X'\eta&\xi'Y'-X'\eta'&\xi''Y'-X'\eta''&\text{products.}
\end{array}
\]

For ease of reference, we denote the sums of products that arise
\begin{center}\begingroup\fontsize{9}{11}\selectfont\setlength{\tabcolsep}{.8pt}
\begin{tabular}{c|c|c|c|c|c}
&&&&&\shortstack{after multiplying\\these factors by}\\\hline
&$*$&$*$&$\delta'\delta D'[\pi'\pi P']$&$\delta''\delta D'[\pi''\pi P']$&$\xi,\quad\eta,\quad\zeta$\\
$[\pi\pi'\pi'']$&$\delta\delta'\delta''\times[\pi\pi'\pi'']$&$\delta\delta'D'[\pi\pi'P']$&$*$&$\delta''\delta'D'[\pi''\pi'P']$&$\xi',\quad\eta',\quad\zeta'$\\
&$*$&$\delta\delta''D'[\pi\pi''P']$&$\delta'\delta''D'[\pi'\pi''P']$&$*$&$\xi'',\quad\eta'',\quad\zeta''$\\
$[\pi P\pi'']$&$\delta D\delta''[\pi P\pi'']$&$\delta DD'[\pi PP']$&$\delta'DD'[\pi'PP']$&$\delta''DD'[\pi''PP']$&$X,\quad Y,\quad Z$\\
$[\pi P'\pi'']$&$\delta D'\delta''[\pi P'\pi'']$&$*$&$*$&$*$&$X',\quad Y',\quad Z'$\\
&$\delta D''\delta''[\pi P''\pi'']$&$\delta D''D'[\pi P''P']$&$\delta'D''D'[\pi'P''P']$&$\delta''D''D'[\pi''P''P']$&$X'',\quad Y'',\quad Z''$
\end{tabular}\endgroup\end{center}
It is clear that the places marked here with $*$ must be filled with $0$, and that all the quantities denoted by bracketed symbols are \emph{given}. Specifically,
\[
\begin{aligned}\relax
[\pi\pi'\pi'']&=\operatorname{tang}\beta\sin(\lambda'-\lambda'')+\operatorname{tang}\beta'\sin(\lambda''-\lambda)+\operatorname{tang}\beta''\sin(\lambda-\lambda')\\
[\pi P\pi'']&=\operatorname{tang}\beta\sin(L-\lambda'')+\operatorname{tang}B\sin(\lambda''-\lambda)+\operatorname{tang}\beta''\sin(\lambda-L)\\
&\text{and so on}
\end{aligned}
\]

There is no need to write down all 16 equations, since they can all be derived from the first in the same way simply by replacing $\beta$ by
\expandafter\def\csname HistoricalNoteText9\endcsname{For brevity I take the liberty of using this readily understood expression. The positive pole of $\mathfrak{X}$ lies on the side of this plane on which the $x$ are counted as positive, and so on.}%
\SourceLeaf{160}%
$\beta'$, $\beta''$, $B$, $B'$, $B''$ and $\lambda$ by $\lambda'$, $\lambda''$, $L$, $L'$, $L''$ when $\pi$ is replaced by $\pi'$, $\pi''$, $P$, $P'$, $P''$, respectively, and so on. At the same time, one sees that the 16 quantities reduce to 12, since
\[
\begin{aligned}
+[\pi P'\pi'']&=-[\pi\pi''P']=+[\pi''\pi P']\\
[\pi\pi'P']&=-[\pi'\pi P']\\
[\pi''\pi'P']&=-[\pi'\pi''P']
\end{aligned}
\]

Furthermore, it is easy to see that the expression $[\pi\pi'\pi'']$, multiplied by the product of the three cosines of the latitudes occurring in it, is six times the volume of a pyramid whose apex lies at the centre of a sphere of radius $1$ and whose three base vertices lie on its surface so as to correspond to the three geocentric positions of $p$. That symbol gives six times the volume with a positive or negative sign according as the three geocentric positions on the sphere lie in the opposite or the same order (sens) as the positive\HistoricalNote{9} poles of the planes $\mathfrak{X}$, $\mathfrak{Y}$, $\mathfrak{Z}$, respectively. The other symbols express entirely analogous quantities.

In this way we obtain the following four equations:
\[
\begin{aligned}
1)\quad &(F+F'')f'\delta'[\pi\pi'\pi'']\\
&=(Ff''-F''f)(D[\pi P\pi'']-D''[\pi P''\pi''])\\
&\qquad+(F'(f+f'')-(F+F'')f')D'[\pi P'\pi'']\\
2)\quad &(F+F'')(f'\delta'[\pi\pi'P']+f''\delta''[\pi\pi''P'])\\
&=(Ff''-F''f)(D[\pi PP']-D''[\pi P''P'])\\
3)\quad &(F+F'')(f\delta[\pi'\pi P']+f''\delta''[\pi'\pi''P'])\\
&=(Ff''-F''f)(D[\pi'PP']-D''[\pi'P''P'])\\
4)\quad &(F+F'')(f\delta[\pi''\pi P']+f'\delta'[\pi''\pi'P'])\\
&=(Ff''-F''f)(D[\pi''PP']-D''[\pi''P''P'])
\end{aligned}
\]

\begin{center}5.\end{center}
We shall now examine these four equations, which are strictly correct, more closely so as to base the first approximation on them. If we regard the intervals between observations as infinitesimal quantities of the first order, then $f$, $f'$, $f''$, $G$, $G'$, $G''$ and all the bracketed quantities are of that same first order, except for $[\pi\pi'\pi'']$, which is of the third order. I omit the proofs and the readily suggested remarks on particular excep%
\SourceLeaf{161}%
tions. If the mutual inclination of the orbits of $p$ and $P$ had to be regarded as a quantity of the first order, all the bracketed quantities would be one order higher. Furthermore,
\[
Ff''-F''f=\frac FG.\frac{f''}{g''}Gg''-\frac{F''}{G''}.\frac fg G''g
\]
or (since $Gg''=G''g$)
\[
=\left(\frac FG.\frac{f''}{g''}-\frac{F''}{G''}.\frac fg\right)Gg''
\]
Now $G-F$ is a quantity of the third order, so $1-\frac FG$ is of the second order, and so on. Consequently $\frac FG.\frac{f''}{g''}-\frac{F''}{G''}.\frac fg$ is also of the second order, and hence $Ff''-F''f$ is of the fourth order (it would even be of the fifth order if $\tau'$ fell midway between $\tau$ and $\tau''$). Thus the right-hand side of the second, third and fourth equations above is of the fifth order, while the first and second parts on the left-hand side are each of the third order. For a first approximation, one can therefore set
\[
\begin{array}{ll}
\text{from 2)}&f'\delta'[\pi\pi'P']=-f''\delta''[\pi\pi''P']\\
\text{from 4)}&f\delta[\pi\pi''P']=-f'\delta'[\pi'\pi''P']
\end{array}
\]
What follows from 3) is identical with these results. To shorten the calculation further, we can also put $\frac fg=\frac{f''}{g''}$ here. Both quantities differ from unity only by a quantity of the second order and differ from each other by the same order (if $\tau'$ falls midway between $\tau$ and $\tau''$, the difference is only of the third order). Since $g:-g':g''=\tau''-\tau':\tau''-\tau:\tau'-\tau$, we obtain
\[
\begin{aligned}
5)\qquad\delta&=\frac gf.\frac{f'}{g'}.\frac{\tau''-\tau}{\tau''-\tau'}.\frac{[\pi'\pi''P']}{[\pi\pi''P']}.\delta'\\
6)\qquad\delta''&=\frac{g''}{f''}.\frac{f'}{g'}.\frac{\tau''-\tau}{\tau'-\tau}.\frac{[\pi\pi'P']}{[\pi\pi''P']}.\delta'
\end{aligned}
\]

These formulas give $\delta$ and $\delta''$ from $\delta'$ correctly up to and including the second order if $\tau'$ lies midway between $\tau$ and $\tau''$, and up to but excluding the second order otherwise. In the latter case one can put $\frac{f'}{g'}=1$, since the difference is only of the second order. In the former case, however, it is worth setting
\[
f'=g'+\tfrac43(f+f'+f'')\quad\text{or}\quad\frac{f'}{g'}=1+\tfrac43\frac{f+f'+f''}{f'}
\]
which will soon be specified more precisely and is accurate to 1 order higher. (It is easy to see that $f+f'+f''$ equals the area of the triangle between the three positions
\SourceLeaf{162}%
\expandafter\def\csname HistoricalNoteText10\endcsname{Abhandlung über die leichteste Methode u.\ s.\ w., p.\ 45.}%
$p$, $p'$, $p''$, so by a familiar approximation it is $=\frac34\times$ the area of the segment bounded by the chord $pp''$ and the arc.) The formulas above also give
\[
\frac{\delta''}{\delta}=\frac{[\pi\pi'P']}{[\pi'\pi''P']}.\frac{\tau''-\tau'}{\tau'-\tau}
\]
which, if one takes the ecliptic for $\mathfrak{Z}$ or sets $B$, $B'$, $B''=0$, reduces to
\[
\frac{\delta''}{\delta}=\frac{\operatorname{tang}\beta\sin(\lambda'-L')-\operatorname{tang}\beta'\sin(\lambda-L')}{\operatorname{tang}\beta'\sin(\lambda''-L')-\operatorname{tang}\beta''\sin(\lambda'-L')}.\frac{\tau''-\tau'}{\tau'-\tau}
\]
that is, to the familiar \textsc{Olbers} formula\HistoricalNote{10}.

\begin{center}6.\end{center}
Having derived convenient approximations from formulas 2.\ 3.\ 4., we now turn to the first in the same way. As is well known,
\[
\begin{aligned}
\frac1r&=\frac1k(1-e\cos(v-\pi))\\
\frac1{r'}&=\frac1k(1-e\cos(v'-\pi))\\
\frac1{r''}&=\frac1k(1-e\cos(v''-\pi))
\end{aligned}
\]

Multiplying by $\sin(v''-v')$, $\sin(v-v'')$, $\sin(v'-v)$ and adding the results gives
\[
\begin{aligned}
\frac{f+f'+f''}{rr'r''}&=\frac1k(\sin(v''-v')+\sin(v-v'')+\sin(v'-v))\\
&=-\frac4k\sin\tfrac12(v''-v')\sin\tfrac12(v-v'')\sin\tfrac12(v'-v)
\end{aligned}
\]
or
\[
\frac{f+f'+f''}{f'}=-\frac{2r'}k.\frac{\sin\tfrac12(v''-v')\sin\tfrac12(v'-v)}{\cos\tfrac12(v''-v)}
\]
By a familiar theorem of theoretical astronomy,
\[
\frac{\text{Area swept out}}{\text{Mean motion}}=\frac{a^{\frac32}.\BareRoot{k}}2
\]
Consequently,
\[
k=\frac{4gg''}{a^3(m'-m)(m''-m')}=\frac{4gg''}{A^3(M'-M)(M''-M')}
\]
\SourceLeaf{163}%
\expandafter\def\csname HistoricalNoteText11\endcsname{This really is subtraction, since $[\pi P\pi'']$, $[\pi P''\pi'']$ have the same sign. By contrast, the terms in the coefficients of $Ff''-F''f$ in equations 2.\ 3.\ 4.\ are in effect added. A more detailed investigation would take us too far here.}%
Thus,
\[
\begin{aligned}
\frac{f+f'+f''}{f'}&=-\frac{A^3(M'-M)(M''-M)}{2\cos^2\tfrac12(v''-v)}.\frac1{r'^3}.\frac{r'r'}{rr''}\\&\quad.\frac{r'r''\sin\tfrac12(v''-v')}{g}.\frac{rr'\sin\tfrac12(v'-v)}{g''}\\
\frac{f+f'+f''}{g'}&=-\frac{r'r''\sin\tfrac12(v''-v')}{g}.\frac{rr''\sin\tfrac12(v''-v)}{g'}.\frac{rr'\sin\tfrac12(v'-v)}{g''}\\&\quad.\frac{A^3(M'-M)(M''-M)}{2r'^3}.\frac{r'r'}{rr''}
\end{aligned}
\]
It follows readily from this that the quantities
\[
\frac1{\cos\tfrac12(v''-v)},\qquad\frac{r'r''\sin\tfrac12(v''-v')}{g},\qquad\frac{rr'\sin\tfrac12(v'-v)}{g'}
\]
and also, if either $\tau'$ falls midway between $\tau$ and $\tau''$ or the departure of the orbit of $p$ from a circle can be regarded as being of the first order, $\frac{r'r'}{rr''}$ differ from unity only by quantities of the second order. One may therefore put, approximately,
\[
\frac{f+f'+f''}{f'}=-\frac{A^3}{2r'^3}(M'-M)(M''-M')
\]
Similarly, one has approximately
\[
\frac{F+F'+F''}{F'}=-\frac{A^3}{2R'^3}(M'-M)(M''-M')
\]

The latter quantity may also be calculated exactly if desired, since everything needed for it is given. Both are of the second order and are determined up to but excluding the fourth order. We therefore have
\[
\begin{aligned}
F'(f+f'')-(F+F'')f'&=F'(f+f'+f'')-(F+F'+F'')f'\\
&=F'f'\tfrac12A^3(M'-M)(M''-M')\left(\frac1{R'^3}-\frac1{r'^3}\right)
\end{aligned}
\]
This quantity is of the fourth order and is determined up to but excluding the sixth. In equation 1 above, the left-hand side is of the fifth order. On the right-hand side, the first term is of the sixth or seventh order, since $Ff''-F''f$ is of the fourth or fifth order and $D[\pi P\pi'']-D''[\pi P''\pi'']$ is of the second\HistoricalNote{11}. The second term is of the fifth order. We therefore drop the first term and obtain
\[
7)\qquad\frac{[\pi\pi'\pi'']}{[\pi P'\pi'']}.\frac2{A^3(M'-M)(M''-M')}=\left(\frac1{R'^3}-\frac1{r'^3}\right)\frac R{\delta'}
\]
\SourceLeaf{164}%
with accuracy up to but excluding the second order if $\tau'$ falls midway between $\tau$ and $\tau''$. Otherwise the error is only of the first order. This formula, which takes the following form if we take the ecliptic for $\mathfrak{Z}$, is the most important part of the whole method and its first foundation.
\[
\begin{aligned}
\left\{1-\left(\frac{R'}{r'}\right)^3\right\}.\frac{R'}{\delta'}&=\frac2{A^3.(M'-M)(M''-M')}\\
&\quad.\frac{\operatorname{tg}\beta'\sin(\lambda''-\lambda)-\operatorname{tg}\beta\sin(\lambda''-\lambda')-\operatorname{tg}\beta''\sin(\lambda'-\lambda)}{\operatorname{tg}\beta\sin(L'-\lambda'')-\operatorname{tg}\beta''\sin(L'-\lambda)}
\end{aligned}
\]
$L'$ is the longitude of the Sun $+180^0$.

Since the right-hand side is given, one can see that combining this equation with the following one
\[
\frac{\frac{R'}{\delta'}}{\frac{R'}{r'}}=\BareRoot{(1+\operatorname{tang}\beta'^2+\frac{R'R'}{\delta'\delta'}+2\frac{R'}{\delta'}\cos(\lambda'-L'))}
\]
will make it easy to find $r'$. The indirect method is by far the most convenient here. A few trials, for which suitable rules are easy to give, lead very quickly to the result. One can also always see whether there is more than one value of $r'$ and hence more than one orbit that represents the observations, as can indeed sometimes happen. —

One should also note that the longitudes ought properly to be measured from a fixed point in place of the moving equinoctial points. In practice, however, the difference is immaterial. If time is expressed in days, one has
\[
\log(M'-M)(M''-M')=\log(\tau'-\tau)+\log(\tau''-\tau')+6.4711352(-10)
\]
(The quantities $M$ and so on must be expressed in parts of the radius in place of degrees.)

Once $\delta'$ and $r'$ are known, one can also determine $\frac{f+f'+f''}{f'}$, and hence $\delta$ and $\delta''$. Examination of formula 7) also yields further interesting consequences, which must be passed over here.

\begin{center}Second stage.\\\emph{Finding approximate orbital elements.}\end{center}
\begin{center}7.\end{center}
We now leave out the middle observation at time $\tau'$ altogether and use the distances $\delta$, $\delta''$ found approximately in the previous stage in its place. %
\SourceLeaf{165}%
Clearly, from these we can now derive the heliocentric longitudes, latitudes and distances, and from these in turn the longitude of $\AscendingNode$, the inclination of the orbit and the longitudes in the orbit. This leaves the task of determining
\begin{center}\begin{tabular}{ll}
from two longitudes in the orbit $\quad.\quad.$&$v,\ v''$\\
\quad the distances from the Sun $\quad.$&$r,\ r''$\\
\quad the corresponding times $\quad.\quad.\quad.$&$\tau,\ \tau''$
\end{tabular}\end{center}
the remaining elements, namely $a$, $e$, $\pi$ and the epoch. Since the relations of these quantities to the given ones are transcendental, one must again turn to indirect methods. We shall consider three here.

\par\medskip\emph{First method.}\par
Assume $\pi$ is given.

Set
\[
\frac{r''+r}{r''-r}\operatorname{tang}\tfrac12(v''-v)=\operatorname{tang}\zeta
\]
Then
\[
e=\frac{\cos\zeta}{\cos\tfrac12(v''-v)\cos[\pi-\tfrac12(v+v'')-\zeta]}
\]
It is then best to calculate $k$ in two ways
\[
k=r[1-e\cos(v-\pi)]=r''[1-e\cos(v''-\pi)]
\]
This also serves to check the calculation. If we set $e=\sin\varphi$, then $a=\frac{k}{\cos\varphi^2}$. The eccentric and mean anomalies and longitudes can then be calculated from the true anomalies by the usual methods or, more conveniently, by indirect methods. From these and the mean motion given by $a$, one obtains two determinations of the mean longitude for any chosen epoch. If they agree, one has found the correct value of $\pi$. Otherwise the calculation must be repeated with a slightly different value of $\pi$ and the true value found by interpolation. ‘It is advisable here to calculate the remaining elements afresh from the corrected values of $\pi$, \emph{without} interpolation, and to continue until the two values for the epoch agree completely.’
\SourceLeaf{166}%
\par\medskip\emph{Second method.}\par
Assume $e$ is given.

The calculation here is exactly the same, except that because $\pi$ must be found from the equation
\[
\cos[\pi-\tfrac12(v+v'')-\zeta]=\frac{\cos\zeta}{e\cos\tfrac12(v''-v)}
\]
one must already know the true value approximately, since two different values correspond to the cosine.

Method I is preferable to II. In general, both methods are suitable only when the arc already traversed is very large and the elements are already nearly known. For the first approximations from a short series of observations, one always uses the following method.

\Needspace{6\baselineskip}\par\medskip\emph{Third method.}\par
Assume $k$ is given.

Here one has
\[
\begin{aligned}
\frac{\frac1{r''}-\frac1r}{2\sin\tfrac12(v''-v)}&=\frac ek\sin(\tfrac12(v+v'')-\pi)\\
\frac{\frac2k-\frac1{r''}-\frac1r}{2\cos\tfrac12(v''-v)}&=\frac ek\cos(\tfrac12(v+v'')-\pi)
\end{aligned}
\]
Division thus gives $\operatorname{tang}[\tfrac12(v+v'')-\pi]$, from this $\pi$, and then $e$ from either equation. The rest is exactly as in the preceding methods.

The advantage of the third method is that one can find a close approximation for $k$ immediately when the arc $v''-v$ is not too large. Indeed, the sector between the two radius vectors is given by
\[
g'=\frac{a^{\frac32}\BareRoot{k}}2(m''-m)=\tfrac12A^{\frac32}(M''-M)\BareRoot{k}
\]
Now $2g'=\int rr\,\mathrm{d}w$ from $w=v$ to $w=v''$.

By the familiar \textsc{Cotes} formula for approximate integration, $\int\varphi w.\mathrm{d}w$ from $w=v$ to $w=v''$ is
\[
=(\tfrac12\varphi v+\tfrac12\varphi v'')(v''-v)
\]
\SourceLeaf{167}%
and more accurately
\[
=(\tfrac16\varphi v+\tfrac23\varphi\tfrac12(v+v'')+\tfrac16\varphi v'')(v''-v)
\]
more accurately still
\[
=(\tfrac18\varphi v+\tfrac38\varphi(\tfrac23v+\tfrac13v'')+\tfrac38\varphi(\tfrac13v+\tfrac23v'')+\tfrac18\varphi v'')(v''-v)\quad\text{and so on}
\]
It is always enough, however, to stop at the first two.

The first therefore gives
\[
2g'=\tfrac12(rr+r''r'')(v''-v)\quad\text{and}\quad\BareRoot{k}=\frac{\tfrac12(rr+r''r'')}{A^{\frac32}}.\frac{v''-v}{M''-M}
\]
One usually sets $A$ $=1$. With $v''-v$ and $M''-M$ expressed in seconds of arc, one has $\log(M''-M)=\log(\tau''-\tau)+3.5500073$. To shorten the calculation one puts $\frac{r'}r=\operatorname{tang}(45^0\pm\psi)$, giving $\tfrac12(rr+r''r'')=\frac1{\cos2\psi}$.

To use the second integration formula, set the radius vector corresponding to longitude $\tfrac12(v''+v)$ to $=r^*$. Then
\[
\frac1{r^*}=\tfrac12\left(\frac1r+\frac1{r''}\right)+\left[\tfrac12\left(\frac1r+\frac1{r''}\right)-\frac1k\right]\frac{2\sin\tfrac14(v''-v)^2}{\cos\tfrac12(v''-v)}
\]
From this one can determine $r^*$ using the first value of $k$. One then has
\[
2g'=\tfrac16(rr+r''r'')+\tfrac23r^*r^*
\]
so the new value of $k$ is
\[
=k\left\{1+\frac{\tfrac23(r^*r^*-\tfrac12(rr+r''r''))}{rr+r''r''}\right\}^2
\]

In practice it is usually accurate enough and more convenient to find the logarithm of the new value of $k$ by increasing the logarithm of the first value by $\tfrac43\log\frac{r^*r^*}{\tfrac12(rr+r''r'')}$. If one uses this new value of $k$ to determine $\frac1{r^*}$ more accurately once again by the formula above and then corrects $k$ once more, the two resulting determinations of the epoch will nearly always agree so closely that no fresh assumption will be needed. For $\text{\AstroSymbols ⚳}$ and $\text{\AstroSymbols ⚴}$, even though $\tau''-\tau$ was 41 and 42 days, the two determinations always agreed to within a couple of hundredths of a second.
\SourceLeaf{168}%
\Needspace{8\baselineskip}
\begin{center}\emph{Refinement methods.}\end{center}
\begin{center}8.\end{center}
If one calculates the position at time $\tau'$ from the approximate elements found by the preceding methods and finds it agrees with the observation, the work is finished. The agreement will usually be very close (in my calculations the difference was often only a few seconds of arc), though it will rarely be perfect, partly because some of the underlying assumptions are approximate and partly \emph{because even the solar positions used in the calculation depart from an ellipse and include the small perturbations}. One could use the approximate elements to determine very closely the values of the small higher-order quantities omitted above and refine the above formulas and the values of $\delta$, $\delta''$ accordingly. In my view, however, this calculation would be far more laborious than one of the following methods.

The easiest method of all for a first refinement is the following. It only occurred to me while working on $\text{\AstroSymbols ⚴}$, and I applied it there with the greatest success, since the intervals were still short.

Suppose that, from the approximate elements \emph{found in the way described above}, the position calculated for time $\tau$ is in longitude $=\lambda'+\mathfrak{L}$ and in latitude $=\beta'+\mathfrak{B}$, whereas the observed values are $\lambda'$ and $\beta'$. All the small inaccuracies in the assumptions thus conspire to make the longitude too large by $\mathfrak{L}$ and the latitude by $\mathfrak{B}$. One should then calculate the orbit all over again and \emph{in exactly the same way}, using the following observations as the basis
\[
\begin{array}{ccc}
\lambda,&\lambda'-\mathfrak{L},&\lambda''\\
\beta,&\beta'-\mathfrak{B},&\beta''
\end{array}
\]
The position calculated from the resulting new elements will differ so little from $\lambda'$, $\beta'$ (in my experience only by fractions of a second of arc) that no further refinement is needed.

\begin{center}9.\end{center}
The procedure just described applies only when the orbit is to be determined from the \emph{same} observations that were used for the first approx%
\SourceLeaf{169}%
\expandafter\def\csname HistoricalNoteText12\endcsname{See \textit{Theoria motus corporum coelest.}, art.\ 83.}%
imation. If one later wishes to refine the elements using observations that are entirely or partly different, then of the various procedures I have tried I have found the following two methods the most useful.

I. One calculates the heliocentric positions from the two outer geocentric positions under 3 hypotheses. First one assumes the approximate distances for these observations, then changes first one distance and then the other slightly. Using the elements found under each of the 3 hypotheses, one calculates the position for the middle observation and compares it with the observed position. Interpolation then yields the corrected distances and, if desired, the corrected elements. It is better, however, to take the trouble to calculate the elements separately from the new distances, especially while the changes in the elements are still very large.

II\textsuperscript{a}. The procedure is exactly the same, with the sole difference that one uses the approximate determinations of the inclination and $\AscendingNode$ in place of the approximate distances for the outer observations, and varies each of these slightly.

II\textsuperscript{b}. One calculates the heliocentric positions from all three geocentric positions, using partly the approximate determinations of the inclination and $\AscendingNode$ and partly slightly altered determinations. From the two outer heliocentric positions one calculates the elements, and from these the middle heliocentric position. One compares this with the position derived from the observed geocentric position, then seeks the improved inclination and $\AscendingNode$ by interpolation, and so on.

In II\textsuperscript{b} one could also use familiar formulas to determine the ellipse from the \emph{three heliocentric} positions without using the times, calculate the 2 intervals from the dimensions of the ellipse and compare them with the actual intervals, then find the corrected inclination and $\AscendingNode$ by interpolation just as before. My experience has forced me to reject this procedure, however. It would take repeated operations and much more effort to represent the observations accurately in this way. To examine the reasons in detail here would take us too far\HistoricalNote{12}. I will only remark that in this way one brings back the second differen%
\SourceLeaf{170}%
tials that were dispensed with precisely by the devices used in Articles 5.\ and 6. These delicate second differentials can be distorted \emph{enormously} by even a modest change in the inclination and $\AscendingNode$, especially when the eccentricity is small. A change of a few minutes of arc in $\AscendingNode$ or the inclination can easily produce an ellipse with scarcely any resemblance to the preceding one. Interpolation, understandably, can then no longer be trusted. This does not happen with our methods, which are always based on only two observations. \emph{Sapienti sat.} — From repeated experience I find method I the most suitable and most general of all.

All these methods apply only while the arc is still of moderate size. Once one has observations from 1 year or more, other methods are needed, which I cannot go into at length here. In that case it is generally inadvisable to base the elements on three complete observations. It is far more appropriate to use four longitudes and two latitudes. — If the observations already span \emph{several} years and the elements have already been determined \emph{apart from small details}, I consider the use of differential corrections, which can be based on any number of observations, the best means.
\par\nobreak\vskip6pt\nobreak\nointerlineskip\hbox to\linewidth{\hfil\rule{2.65cm}{0.4pt}\hfil}
